\chapter{Learning from Examples and from Instructions}\label{ch:human}

People learn a task from examples of it and from being told what it is. This thesis has put a price on the second in units of the first. For the Bayes-optimal predictor, a description is worth a number of demonstrations set by its specificity and its reliability (Chapters~\ref{ch:specificity} and~\ref{ch:reliability}); networks trained on the task distribution come close to that price for reliable descriptions, and pretrained language models answering directly do not (Chapters~\ref{ch:networks} and~\ref{ch:llm}). Psychology has long studied how people learn by each route and how they weigh what they are told, under the names of observational and instructed learning, trust in testimony, and implicit and explicit learning. This chapter places the thesis's results beside that work. Section~\ref{sec:hl-scope} states what kind of comparison this is. Sections~\ref{sec:hl-worth} to~\ref{sec:hl-deliberate} take three questions in turn: what an instruction is worth, what it costs when it is wrong, and how a stated quantity comes to be weighed against experience. Section~\ref{sec:hl-limits} sets out where the parallels fail, and Section~\ref{sec:hl-test} proposes the experiment that would replace them with a measurement.

\section{The Level of the Comparison}\label{sec:hl-scope}
The theorems of Chapters~\ref{ch:specificity} and~\ref{ch:reliability} describe an ideal learner: the predictor that is Bayes-optimal for the task distribution. A comparison with people is therefore a comparison at what \citet{marr1982} called the computational level, the level of the problem being solved and of its ideal solution, not the level of the algorithm or the brain that solves it. A Bayesian account at this level does not assume that people represent a posterior or apply Bayes' rule \citep{griffiths2010}; it says what an optimal use of the information would look like, so that a learner's departures from it can be measured. Meta-trained networks are a way to obtain such ideal learners when the posterior is hard to compute, and \citet{binz2024meta} argue that they can serve as rational models of human learning for that reason. Our networks of Chapter~\ref{ch:networks} are an instance of that program, and the Bayes-optimal predictor is the common reference against which the networks, the language models, and people can each be placed.

Three limits follow. First, this chapter reports no new data from people; every human finding comes from published studies, on tasks that differ from ours, and every statement about how people would behave on our task is a conjecture. Second, a match between a model and people in behaviour is not a match in mechanism; where both depart from the ideal in the same direction, we first ask whether a simpler property of the task or of the information given explains it, as \citet{frank2023} urges for comparisons of this kind. Third, words that describe people are not transferred to models. We describe the networks and language models only by what they predict and how their predictions depend on the prompt. Where the chapter speaks of a model's ``trust'' in a description, it means the weight its predictions put on the description, defined operationally as in Chapters~\ref{ch:networks} and~\ref{ch:llm}; we do not attribute belief, understanding, or awareness to it \citep{shiffrin2023}, and in particular we do not describe any model process as conscious or unconscious \citep{butlin2023}. On terms: \emph{description} remains the thesis's word for its object, while \emph{instruction}, \emph{statement}, \emph{advice}, and \emph{testimony} are the human studies' own words for information about a task or a source that is given in words rather than shown.

Parts of the mapping we draw already exist. \citet{russin2025} relate in-context and in-weights learning in Transformers to the explicit and implicit learning systems of human dual-system theories, and \citet{hagendorff2023} and \citet{lampinen2024content} read a language model's direct answer and its chain of thought as analogues of fast and deliberate human reasoning. The thesis adds a common unit. The effective sample size gives the value of a stated description in demonstrations, so the question of how people and models weigh what they are told against what they are shown can be asked with a number and an optimum.

\section{What an Instruction Is Worth}\label{sec:hl-worth}

\paragraph{Two routes to a task.} \citet{bandura1961} showed that children acquire new actions by watching an adult perform them, without reward to the adult or to the child, and \citet{bandura1977} built a theory of social learning on the observation that learning only from the consequences of one's own actions would be slow and hazardous. The question of how many demonstrations a description is worth asks for the exchange rate between two routes to the same competence. The idea of measuring prior information in units of hypothetical data is old in statistics. \citet{keynes1921} separated the \emph{weight} of an argument, the amount of relevant evidence behind it, from the probability it supports, and \citet{good1950} proposed calibrating a prior by asking what imaginary results would change one's mind. The effective sample size of Section~\ref{sec:bg-ess} gives Keynes's question one cardinal answer for a particular model, and runs Good's device in reverse: it converts a prior into the data it is worth.

\paragraph{A reliable instruction replaces much experience.} People given instructions for a new task often perform it correctly on the first trial, where monkeys need about twenty trials to reach the same accuracy and early neural network models thousands \citep{cole2013}. When people are told the correct reward probabilities of the options in a choice task, they earn more than when they learn them from outcomes, and their fitted learning rate from feedback falls from $0.24$ to nearly zero \citep{li2011}. This is the human counterpart of a reliable, specific description in Chapter~\ref{ch:specificity}, which at $d = 16$ and $a_0 = 10$ is worth $117$ demonstrations when it leaves a thousandth of the prior variance (Proposition~\ref{prop:reliable}). First-trial success from instructions depends on composing rules the person has already practised \citep{cole2013}, and a description is worth something to a learner only if it is expressed in terms the learner already uses. Our networks receive the description in the coordinates of the task, and models trained only on task definitions cannot use examples, or the reverse, unless training included both \citep{wang2022supernat}.

\paragraph{People have an exchange rate too.} \citet{tessler2020} measured how many examples a generic statement is worth to adults. Participants rated the probability that a new member of a novel category had a property, after reading a generic such as ``Feps have white wings'' or after seeing one to four examples of feps with white wings. When the property was named alongside the examples, a generic was worth about two examples chosen by a teacher, or three encountered by accident. The number is far below the $117$ of Chapter~\ref{ch:specificity}, and the theorems suggest one reason. A description's ESS grows with how much it pins down relative to what the learner already expects, and a generic leaves much unstated. Adults accept a novel generic as true when as few as $10\%$ to $30\%$ of a category have the property, yet take one they are told to imply that nearly all members have it \citep{cimpian2010}. The thesis's descriptions state a value and a spread; a generic states neither, so what it is worth depends on how the learner fills in what it leaves out. \citet{sumers2023} approach the question from the other side. In a game where one player teaches another a rule over objects' colour, shape, and fill, teaching by language beat teaching by demonstration, and its advantage grew with the complexity of the rule, a pattern consistent with a description's ESS growing with what it specifies (Proposition~\ref{prop:reliable}). The advantage largely disappeared when the teacher could see which objects were valuable but not their features, while demonstrations were unaffected. Language transmits a rule only through abstractions that teacher and learner share, so a description is worth many examples only when both know the features it must name.

\paragraph{Examples chosen for the learner are worth more.} The same examples support stronger generalization when learners believe a helpful teacher chose them than when they believe the examples were sampled at random \citep{shafto2014}, and \citet{tessler2020} found that a generic was worth fewer examples chosen by a teacher than examples met by accident. The human literature thus distinguishes three channels: incidental observation, demonstrations selected for the learner, and verbal statements. The demonstrations of this thesis are random draws from the task, the first channel. A demonstration chosen to be informative would carry part of what a description carries, and the ESS of a description would then be smaller against such demonstrations; the thesis does not measure this case.

\paragraph{What the parallel supports.} At the computational level, people and the Bayes-optimal predictor agree that a statement can replace much experience and that it helps only a learner who can interpret it, and the human evidence is consistent with its value growing with how much it specifies. The parallel does not say that people value statements at the Bayes-optimal rate; no study we found measures the human exchange rate for a statement with a stated spread, the case our theorems price exactly.

\section{What an Instruction Costs}\label{sec:hl-cost}
Chapter~\ref{ch:reliability} gave a description that may be irrelevant two prices. A learner that hedges pays $p(1-p)(1-r)\,a_0 d$ for not knowing whether the description is relevant when no demonstrations are in hand ($14.4$ at $r = 0.001$ and $p = 0.9$; Proposition~\ref{prop:unreliable}), and demonstrations shrink this price as they show which reading holds. A learner that trusts every description pays little with no demonstrations but keeps predicting near the stated value, so a specific description leaves it worse off than no description from $15$ demonstrations on (Section~\ref{app:trust}). People offer evidence on both sides of this comparison.

\paragraph{People keep following a wrong instruction.} \citet{doll2009} gave participants a false instruction about which stimulus in a probabilistic choice task was most, or least, likely to be correct. In the last block of training, on the pair the instruction concerned, instructed participants chose the better option only $36.5\%$ of the time, against $53\%$ for participants told nothing, and only five of sixteen chose it on at least half of their last ten choices between the pair. The best-fitting model amplified outcomes consistent with the instruction and discounted inconsistent ones; it fit better than a model in which the instruction sets a prior that the outcomes then override, as in a Bayesian learner. Advice given once before a choice task has the same lasting effect \citep{biele2009, biele2011}, and in the simulations of \citet{biele2011} this bias earned more reward than letting advice set only the initial expectations in most of the environments considered, so that it can be a reasonable policy for a learner whose experience says that advice is usually good. Three-year-olds told by an adult that a sticker was under the wrong cup kept searching there over eight trials, finding the sticker on $1.44$ of them, below chance, while the same false information given by an arrow was soon discounted \citep{jaswal2010}. The over-trusted channel is not always the statement, however. Children copy an adult's visibly unnecessary actions on a puzzle box even after being told to do only what they have to do \citep{lyons2007}, and adults copy such actions more faithfully than children \citep{mcguigan2011}; there a demonstration prevails over a verbal warning, though one that did not say which action was unnecessary. Instruction can also cost what it leaves out. Preschoolers shown one function of a toy by a teacher explored less and found fewer of its other functions than children who saw that function produced by accident or were not shown it at all \citep{bonawitz2011}, and in chains of adults improving a physical device, passing on written theories, mostly oversimplified, alongside the device narrowed what later learners tried \citep{derex2019}.

\paragraph{The language models discount a contradicted description.} On this point the language models of Chapter~\ref{ch:llm} behave more like a learner that allows a description to be irrelevant, the mixture of Chapter~\ref{ch:reliability}, than like the participants of \citet{doll2009}. By two to eight demonstrations, depending on the model, they discount a stated value the demonstrations contradict. At four demonstrations, their weight on an irrelevant $c = 36$ value is $0.01$ to $0.33$, against $0.18$ to $0.44$ on a relevant one (Section~\ref{sec:reasoning}). Their errors therefore stay far below those of the Bayes-optimal predictor of Chapter~\ref{ch:llm}, which takes every description as relevant, partly because they discount the contradicted value and partly because their weight on any description is modest. The difference from people may lie in the task more than the learner. In \citet{doll2009} the evidence against the instruction was weak, a choice between options rewarded $60\%$ and $40\%$ of the time that uninstructed participants also barely learned, whereas our demonstrations of a hidden mean contradict an irrelevant stated value within a few numbers. Large language models also accept a claim in their context less often the more confident they are in their own answer and the further the claim departs from it \citep{wu2024clasheval}. Both observations fit a learner that treats a description as one source of evidence among others; neither shows that the models do so in the way a Bayesian learner does.

\paragraph{Trust is learned from a track record.} With no demonstrations, a prompt carries no information about whether its description is relevant (Section~\ref{app:reliability}, Step~1), so a learner's trust must come from outside the prompt. Our networks take it from training. Those trained on reliable descriptions treat every description as relevant, at an implied trust of $0.95$ to $1.02$, and from six demonstrations on are hurt when one test description in ten is irrelevant, while networks trained at $p = 0.9$ keep hedging, at $0.87$ to $0.89$, even when every test description is relevant (Chapter~\ref{ch:networks}). Children's trust develops in the same way. Four-year-olds prefer to learn from an informant who has labelled objects accurately before, while three-year-olds do not \citep{koenig2005}, and four-year-olds track informants' relative rates of error, while three-year-olds prefer the more accurate informant only when one has made no errors at all \citep{pasquini2007}. \citet{shafto2012trust} model this development as a change in children's default belief about how knowledgeable and helpful informants are, a prior learned from experience. Six-year-olds who had seen a teacher demonstrate only one of a toy's several functions explored more broadly when the same teacher later demonstrated a new toy, a hedge of the kind whose price Proposition~\ref{prop:unreliable} gives \citep{gweon2014}.

The philosophy of testimony has asked the same question about where trust comes from. \citet{hume1748} held that our reliance on testimony rests on experience of how often reports have matched the facts, and \citet{reid1764} that it is an innate disposition, strongest in children and restrained by experience. In the terms of this thesis, a learner's default trust is Humean in origin. The networks' trust is set by the frequency with which descriptions fitted tasks in training, and the Bayes-optimal learner's $p$ is itself such a frequency. Only within a prompt does the Bayes-optimal learner behave as Reid described, beginning inclined to the description and held back as contrary demonstrations accumulate. The robust mixture of Chapter~\ref{ch:setup} has the skeleton of the partially reliable source of \citet{bovens2002}, a source that is either fully reliable or no better than a randomizer, extended so that a relevant description has a finite precision. With that extension, the cost of full trust depends on how specific the description is.

\paragraph{Two kinds of vigilance.} \citet{sperber2010} distinguish vigilance towards the source of a message, deciding whom to believe, from vigilance towards its content, checking what is said against what else one knows. The two language-model findings separate them. The models discount a description that the demonstrations contradict, which is vigilance towards content, but no model hedges in proportion to a stated chance that the description is wrong, which is information about the source. Stating a $10\%$ or $50\%$ chance leaves the ESS of gpt-oss-20b and OLMo~2 32B about where it was and raises Phi-4's; Qwen3.5 9B and Gemma~4 31B lower it, but by about as much for $50\%$ as for $10\%$; Gemma~4 12B lowers it for $10\%$ and raises it for $50\%$; the Bayes-optimal ESS falls from $36$ to $22$ and $5.8$ (Section~\ref{sec:reasoning}). People respond to cues about a source's reliability, if weakly. Across $129$ data sets, the weight people give to advice rises from $0.32$ to $0.48$ between advisers presented with cues of low and of high quality, while the advice's actual accuracy, which could be coded for only about a third of the effect sizes, was not a unique predictor \citep{bailey2023}. Five- and six-year-olds use an informant's stated confidence; only adults also check whether that confidence has been warranted, and they do so less under a concurrent load \citep{tenney2011}. People thus respond to what they are told about a source, though no study we found gives a benchmark for how much they should; the models answering directly respond in ways that bear little relation to the size of the stated chance.

\section{Weighing a Stated Quantity}\label{sec:hl-deliberate}
In Chapter~\ref{ch:llm}, the language models answering directly respond to a description's stated spread far less than the Bayes-optimal predictor does. With one demonstration shown, the Bayes-optimal weight on the stated value rises from $0.50$ to $0.80$ to $0.97$ as the stated standard deviation narrows. The weight of gpt-oss-20b, Phi-4, and OLMo~2 32B changes by at most $0.08$; the Gemma models' weight rises by about $0.25$, half the Bayes-optimal rise; Qwen3.5's does not rise (Section~\ref{sec:reasoning}). With room to reason first, gpt-oss-20b puts weights of $0.50$, $0.78$, and $0.95$ on the same values, and GPT-5.6-Luna matches the Bayes-optimal weight to within $0.03$ (Sections~\ref{sec:reasoning} and~\ref{sec:llm-runs}). The networks of Chapter~\ref{ch:networks}, which learn the prior from training, come close to the Bayes-optimal ESS for reliable descriptions without any such step. The contrast between information that arrives as experience and information that arrives as a stated quantity has a long history in the study of people.

\paragraph{Priors learned from experience and priors that are stated.} People who experience a distribution of perturbations over many trials of a reaching task weight what they have learned against noisy feedback roughly as a Bayesian observer would \citep{kording2004}, although a reanalysis of the individual trials questions how close to optimal this is \citep{duffy2022}. Stated base rates fare worse. Told that a group contains $70\%$ or $30\%$ engineers, people judged an uninformative description of one member to be equally likely an engineer under both rates \citep{kahneman1973}, and in a meta-analysis of updating experiments with stated priors, the regression weight on the prior's log odds is $0.60$ rather than the Bayesian $1$ \citep{benjamin2019}, a weight measured, like ours, against its Bayesian value. \citet{koehler1996} concluded that base rates are used, but more when they are learned from experience or expressed as frequencies than when they are stated as probabilities. Restating the same numbers as natural frequencies, the counts that accumulated experience would give, raised the share of Bayesian answers from $16\%$ to $46\%$ \citep{gigerenzer1995}. The format of the same information changes how much weight it receives. The gap between decisions from described and from experienced risks \citep{hertwig2004} arises in large part because small experienced samples differ from the described distribution \citep{fox2006}, so a fair comparison holds the information fixed. The ESS does so by construction, since it compares a description with the demonstrations that would carry the same information to the same learner.

\paragraph{People adjust to a stated spread; some models answering directly barely do.} The closest human experiment to ours shows people a prior distribution explicitly, as a visual density whose width changes from trial to trial, together with a noisy cue \citep{acerbi2014}. Participants' weight on the cue rose with the width of the prior, in the Bayesian direction, though they gave the prior less weight than it deserved. Here four of the six models, answering directly, part company with people, while the Gemma models, whose weight rises by about half the Bayes-optimal amount, look more like the participants of \citet{acerbi2014}. The models share with people a weight that responds to the stated precision less than it should, but their weight is not uniformly too low. With one demonstration shown, all but the Gemma models put more weight on the least specific value than the Bayes-optimal $0.50$, at $0.53$ to $0.82$ (Table~\ref{tab:llm}). A weight that barely responds at all appears in some models and not, in this experiment, in people.

\paragraph{Told but not used.} In the study of implicit learning, people sometimes acquire a stated description of a task without using it to perform the task. \citet{berry1984} gave participants who were learning to control a simulated production system a verbal description of the equation that governed it. The description improved their answers to questions about the system but not their control of it; the description combined with thinking aloud during the task did improve control. \citet{sanchez2013} had participants memorize a repeating sequence before practising a fast sequence-tracking task; they recalled it well but tracked it no better than participants who had never been told it, and the authors suggest that applying explicit knowledge is slower than the task allowed. In both, as in our models, a stated quantity is available to the learner but does not shape its answer until the learner works with it. We do not claim more than this resemblance. Thinking aloud is not a chain of thought, both studies are small, and the human results concern control and timing rather than the weighing of a prior.

\paragraph{What ``conscious'' can and cannot mean here.} In this literature, learning that improves performance without improving what people can report is called implicit, and implicit learning was long described as unconscious \citep{reber1967}. That description is disputed. Critics argue that the tests of awareness were insensitive or probed the wrong knowledge, and that performance often rests on fragments of the examples that people can report \citep{shanks1994, perruchet1990, vadillo2016}. Dual-process accounts of reasoning now define the deliberate type of processing by its use of working memory to hold a representation apart from the default response, and treat consciousness only as a typical correlate whose definition is problematic \citep{evans2013}. That functional definition is the one that transfers. Reasoning text gives a model a place to hold the stated prior and combine it with the demonstrations before answering, a role akin to working memory, and in gpt-oss-20b and GPT-5.6-Luna, the answers given after reasoning put close to the Bayes-optimal weight on the stated value. We did not inspect the reasoning text. One reading, a hypothesis rather than a finding, is that it applies the textbook formula for a normal mean (Section~\ref{sec:bg-ess}), so that the stated spread enters the answer once it is put into a computation. The parallel says nothing about awareness, which the analysis of \citet{butlin2023} suggests no AI system of 2023 had, and nothing about a separate deliberate system, since one network produces both kinds of answer \citep{lampinen2024content}.

\paragraph{The pattern, narrowed.} It would be tempting to conclude that experience is used well without deliberation and statements only with it. Neither half holds in general. People learning from small samples of experience underweight rare events \citep{hertwig2004} and can learn local fragments of examples rather than their rule \citep{perruchet1990}. People apply a simple stated rule quickly, though similarity to the examples they saw still pulls at their answers \citep{hahn2010}. Deliberation can also hurt. Chain-of-thought lowers language models' accuracy on tasks where thinking lowers people's, such as learning an artificial grammar from examples \citep{liu2025mind}. The thesis supports a narrower claim, about the models only: a deliberate step helped them weigh a stated spread against data. Whether the same holds for people is open. Thinking aloud let people use a stated equation \citep{berry1984}, and adults check a speaker's calibration only when not under load \citep{tenney2011}, both consistent with it. But on base-rate problems, the canonical task of weighing a stated prior against case evidence, people who gave a wrong first answer, under a deadline in some of the experiments, seldom corrected it after deliberating, and most correct final answers had already been given fast \citep{bago2017}. Either way, the thesis contrasts two formats, a statement and data, within one learning system; in the terms of \citet{russin2025}, both our demonstrations and our descriptions are learned in context.

\section{Where the Parallels Fail}\label{sec:hl-limits}
\paragraph{Different measures.} The thesis measures a learner by its squared error and converts the difference into demonstrations. Human studies measure accuracy, choice rates, or the weight of advice, and a mean weight can hide a mixture of responses. People revising an estimate after advice often either keep their own or adopt the adviser's, rather than averaging \citep{soll2009}, and a regression weight across tasks, ours included, cannot tell such a mixture from uniform partial weighting. The human experiments also differ from ours in how information is given. \citet{acerbi2014} show a prior as a picture over many trials with feedback; we state it once in a sentence and ask for one answer. We found no study that states a prior's spread to people in words and measures the weight they give it, the exact counterpart of Chapter~\ref{ch:llm}.

\paragraph{Different kinds of instruction.} Human instruction is communication. A teacher chooses what to say for a learner, and the learner reasons about why it was said \citep{shafto2014, sperber2010}. Our description is a statement of a task parameter whose relation to the task is fixed by a probability model, and our demonstrations are random draws, not examples chosen to teach. The pragmatic inferences that make a teacher's omission informative \citep{bonawitz2011, gweon2014} have no place in our setting. This is the strongest objection to the comparison, and we grant it. The comparison still holds for the information a statement carries. Whatever else a statement does, it changes what the learner knows about the task, and the ESS prices that change.

\paragraph{Different histories.} The language models were trained on text written by people, so any human-like pattern in their use of a description may have been learned from that text rather than arrived at independently, and our experiments cannot separate the two. The networks of Chapter~\ref{ch:networks}, by contrast, learned from a known distribution, which is what makes their trust interpretable and also what makes them a poor model of a learner whose experience of instructions is varied and social.

\section{A Test with People}\label{sec:hl-test}
The parallels above can be turned into a measurement. The prompts of Chapter~\ref{ch:llm} are a task people can do: estimate a hidden mean from a few numbers, given a sentence that states a value and its spread and, in some conditions, a chance that the sentence is wrong. For lay participants the description would use a wording that does not assume statistical training, such as ``Its mean is about $m$, give or take $\tau$'' (Section~\ref{sec:llm-stated}), after a check that participants understand it. Each participant would answer many items with fresh tasks and no feedback, with the number of demonstrations kept small and the stated spread varied between participants, so that, like the models answering each prompt on its own, no participant can compare descriptions of different spreads. Because people average numbers imprecisely, an ESS read from their own no-description error would mix the description's value with their arithmetic, as it does for the models (Section~\ref{sec:llm-task}), so the primary measure would be the weight on the stated value. Its Bayes-optimal reference would be adjusted for each participant's error in averaging, estimated from the no-description items, since a learner who averages imprecisely should rely more on a stated value. The studies reviewed here suggest two predictions and leave one question open. People's weight would rise with the stated precision, in the Bayesian direction, while the stated value receives less weight overall than the Bayes-optimal predictor gives it, as in \citet{acerbi2014}. Their weight would fall with a stated chance of being wrong, but by less than the stated chance warrants; the advice-taking evidence for this is weaker, since there cues of high quality raise the weight more clearly than cues of low quality lower it \citep{bailey2023}. Whether asking them to work through the numbers before answering brings them closer to the Bayes-optimal weight, as reasoning did for the models, is open \citep{berry1984, bago2017}. Each prediction is untested, and such a study would require research ethics approval, which the work in this thesis did not require.
